% Modernised reading — fonds Alexandre Grothendieck, shelfmark 114, pages 1-5.
%
% This is an INTERPRETATION, not a transcription. Everything the critical
% apparatus carried moves into footnotes. In any dispute of substance the
% transcription governs: whoever cites Grothendieck must cite batch-01.fr.tex.
%
% Pass: Opus 5 (claude-opus-5), 2026-08-17 — first pass, unchecked against the
% pages by a human. The folder was read whole; it holds one batch, pages 1-5.
%
% Two places where this reading departs from the page, both flagged in
% footnotes at the point of departure:
%
% — The « NB » of page 4 asserts that an aspheric functor f induces an f* that
%   is conservative and faithful, and glosses this as « f dominant ». Read as a
%   statement about isomorphisms of presheaves that is false, and a two-object
%   counterexample is given. What dominance does characterise is exactly the
%   bracketed formula the page displays two lines below, which is proved here.
%   The page's own working around that NB is struck through and restarted twice,
%   so the edition is not correcting a settled claim but reporting where a
%   working note was going wrong — and the page reaches the right criterion by
%   its last line.
%
% — Proposition 4's proof turns on the sandwich W_infinity ⊂ W ⊂ W_0. The left
%   inclusion is Grothendieck's own conjecture on the minimal basic localizer,
%   proved by Cisinski twenty years later; the reading says so rather than
%   letting the proposition look self-contained.
\documentclass[11pt,a4paper]{article}
% Le préambule commun à toutes les transcriptions du fonds Grothendieck.
%
% Il tient en une idée : **l'incertitude se compose, elle ne se résout pas.**
% Une transcription qui lisse un mot douteux a détruit la seule chose pour
% laquelle on la faisait — un lecteur doit pouvoir savoir, à chaque endroit,
% ce qui était sur le feuillet et ce qui a été ajouté. Les six macros qui
% suivent sont donc l'appareil critique, et non de la décoration.
%
% Les mêmes commandes sont reconnues par `scripts/render.mjs`, qui produit la
% vue de lecture du site. Une macro ajoutée ici doit l'être aussi là-bas,
% faute de quoi le rendu échouera bruyamment — ce qui est voulu : un rendu
% silencieusement faux est pire qu'un rendu absent.

\ProvidesPackage{grothendieck}[2026/08/08 Transcriptions du fonds Grothendieck]

\usepackage{amsmath,amssymb,amsthm}
\usepackage{mathrsfs}
\usepackage{stmaryrd}
% Les diagrammes commutatifs sont partout dans ce fonds : c'est la langue
% même de la géométrie algébrique de Grothendieck, et une transcription qui
% les rendrait en prose ne transcrirait rien.
\usepackage{tikz-cd}
% Français par défaut — c'est la langue des feuillets — mais l'anglais est
% chargé aussi : la traduction et le résumé sont des documents anglais, et la
% typographie française (espaces avant les deux-points) y serait une faute.
% Ces documents basculent par \selectlanguage{english} en tête de corps.
\usepackage[english,french]{babel}
\usepackage{fontspec}
\usepackage{geometry}
\geometry{a4paper,margin=25mm,marginparwidth=28mm}
\usepackage{marginnote}
\usepackage{xcolor}
% ulem, not soul. `soul`'s \st and \ul are built on a character-by-character
% scan that XeLaTeX's font machinery breaks: the document fails with an
% undefined control sequence rather than degrading. ulem does the same two jobs
% and is Unicode-engine safe. `normalem` stops it hijacking \emph, which the
% transcriptions use for Grothendieck's own emphasis.
\usepackage[normalem]{ulem}
\usepackage{hyperref}
\hypersetup{colorlinks=true,linkcolor=black,urlcolor=black,pdfborder={0 0 0}}

% --- Métadonnées du lot -----------------------------------------------------
% Reprises telles quelles par le script de rendu, qui les affiche en tête de
% la vue de lecture. Elles fixent ce que le fichier prétend couvrir : sans
% elles, une transcription flotte sans point d'attache dans un fonds de seize
% mille feuillets.

\newcommand{\folder}[1]{\gdef\grothFolder{#1}}
\newcommand{\batch}[1]{\gdef\grothBatch{#1}}
\newcommand{\leaves}[2]{\gdef\grothFirst{#1}\gdef\grothLast{#2}}
\newcommand{\foldertitle}[1]{\gdef\grothFoldertitle{#1}}
\gdef\grothFolder{?}\gdef\grothBatch{?}\gdef\grothFirst{?}\gdef\grothLast{?}\gdef\grothFoldertitle{}

% --- Le résumé --------------------------------------------------------------
%
% Il ouvre la lecture modernisée, sous le titre « Résumé », et il est composé
% en petit. Ce n'est pas un ornement : c'est le seul passage du document qui
% ne soit pas des mathématiques, et un lecteur doit pouvoir le reconnaître
% avant de l'avoir lu — puis le sauter s'il connaît déjà le sujet, ou s'y
% arrêter s'il ne le connaît pas. Le filet qui le suit marque la reprise du
% corps.

\newenvironment{resume}
  {\small\setlength{\parskip}{0.45em}}
  {\par\medskip\hrule\medskip}

% --- Mots-clés ---------------------------------------------------------------
%
% En anglais, en clôture du résumé de la lecture modernisée : le vocabulaire
% moderne sous lequel chercher ce que les feuillets construisent. Le manifeste
% du site (`npm run manifest`) les extrait pour étiqueter la cote entière —
% cette ligne est la source unique des tags, il n'y a pas de fichier à côté.
\newcommand{\keywords}[1]{\par\smallskip\noindent
  {\footnotesize\textsc{Keywords} --- \textit{#1}}\par}

% --- L'appareil critique ----------------------------------------------------

% Le numéro de feuillet, en marge. C'est l'ancre qui permet de régler un
% désaccord sur une formule en regardant *un* feuillet plutôt qu'une chemise
% de six cents. Le site s'en sert aussi pour tourner les pages du fac-similé
% quand on fait défiler la transcription.
\newcommand{\leaf}[1]{%
  \marginnote{\footnotesize\color{gray!70}\textsf{#1}}%
  \label{leaf:#1}\ignorespaces}

% Illisible : on ne devine pas. Un mot inventé qui se lit comme les autres est
% le pire résultat possible.
\newcommand{\ill}{\textcolor{red!60!black}{[\,\ldots\,]}}

% Lecture douteuse : le mot est donné, et signalé comme incertain.
\newcommand{\uncertain}[1]{\uline{#1}}

% Ajout de l'éditeur — une lettre manquante, un mot sous-entendu.
\newcommand{\add}[1]{\textcolor{blue!50!black}{[#1]}}

% Biffé par Grothendieck lui-même. Ce qu'il a rayé fait partie du document :
% une rature dit souvent où la pensée a changé de direction.
\newcommand{\struck}[1]{\sout{#1}}

% Note du transcripteur, sur l'état matériel du feuillet.
\newcommand{\note}[1]{\par\smallskip{\small\color{gray!60!black}\textit{[#1]}}\par\smallskip}

% Note marginale de Grothendieck, distincte du corps du texte.
\newcommand{\marginal}[1]{\par\smallskip\noindent\hspace{1em}%
  {\small\color{gray!60!black}#1}\par\smallskip}

% --- Titre ------------------------------------------------------------------

\AtBeginDocument{%
  \begin{center}
    {\large\bfseries Fonds Alexandre Grothendieck --- cote n\textsuperscript{o}~\grothFolder}\\[2pt]
    {\itshape\grothFoldertitle}\\[4pt]
    {\small feuillets \grothFirst--\grothLast\ (lot \grothBatch)}\\[2pt]
    {\footnotesize\color{gray!60!black}Transcription établie d'après le fac-similé numérisé
     par l'Université de Montpellier.\\
     Les lectures incertaines sont soulignées, les illisibles notées [\,\ldots\,],
     les ajouts entre crochets.}
  \end{center}
  \vspace{1em}}

\endinput


\folder{114}
\pages{1}{5}
\dating{[1983]}
\watermark{Édition de démonstration\\interprétation personnelle de l'œuvre}
\foldertitle{Asphéricité : notes manuscrites (s.d.) — lecture modernisée du dossier entier}

\begin{document}

\section*{Résumé}

\begin{resume}
Cinq feuillets, et l'un d'eux n'est pas de la main : c'est une page
dactylographiée, numérotée 350, arrachée à \emph{À la poursuite des champs}. Les
quatre autres tournent autour de ce qu'elle énonce. La chemise porte un seul
mot au crayon : « Asphéricité ».

Le mot demande qu'on dise d'abord de quoi il s'agit. On sait depuis Poincaré
attacher à un espace des invariants algébriques, et l'on sait que des espaces
d'aspect très différent peuvent porter exactement les mêmes : on dit qu'ils ont
le même \emph{type d'homotopie}. Grothendieck poursuit ici l'idée que le type
d'homotopie ne tient ni à l'espace ni aux triangles dont on le recolle, mais se
laisse porter par des \emph{catégories} — et qu'il faut donc savoir lesquelles.
Une petite catégorie $A$ dont les préfaisceaux modélisent l'homotopie est une
\emph{catégorie test}. Un objet est \emph{asphérique} lorsqu'il est, du point de
vue de l'homotopie, aussi simple qu'un point : sans trous, sans composantes
multiples, contractile au sens faible.

Un objet contractile ne se reconnaît pas à l'œil, et surtout il ne se reconnaît
pas indépendamment de ce qu'on a décidé d'appeler « équivalence ». C'est là
qu'est la difficulté de ces pages. Grothendieck ne travaille pas avec une seule
notion d'équivalence mais avec une classe de notions possibles, les
\emph{localisateurs fondamentaux}, chacune donnant sa propre asphéricité. Une
propriété qui dépendrait du choix ne serait pas une propriété de l'objet.

Et c'est ce dont la page 350 s'occupe, dans un paragraphe qui vaut d'être lu
pour lui-même. Grothendieck s'y arrête net : « après une minute de réflexion, je
me sens très bête » de découvrir qu'il traîne depuis des mois la notion de
catégorie totalement asphérique « sans avoir pris la peine de m'asseoir et de
regarder » si elle dépend vraiment du choix. Il s'assoit, il regarde, et la
réponse est non — c'est la proposition 4. Mieux : la condition ne fait plus
intervenir l'homotopie du tout, seulement le comptage des composantes connexes.
Une notion qu'on croyait suspendue à un choix se révèle élémentaire, et
l'aveu de n'avoir pas vérifié plus tôt est laissé sur la page.

Les deux derniers feuillets attaquent autre chose, et sans succès. Un
modélisateur $M$ a une partie asphérique $M_{as}$ ; on voudrait qu'elle suffise
à reconnaître les objets de $M$. Grothendieck croit tenir un critère, l'écrit,
le rature, le reprend, construit deux contre-exemples qui montrent que la partie
asphérique peut être vide, ou réduite à un point sans rien engendrer du tout, et
finit par le bon énoncé — il faut que $M_{as}$ soit \emph{génératrice} — puis par
une question qu'il laisse ouverte au bas du feuillet. On voit rarement d'aussi
près une idée en train de ne pas marcher.

\keywords{test category, basic localizer, total asphericity, modelizer, minimal fundamental localizer, generating family}
\end{resume}

\pagerange{1}{1}
\section*{Le programme (page 1)}

Le premier feuillet ne démontre rien : c'est une liste de neuf questions, en
anglais, avec dans la marge de droite le numéro du chapitre où chacune devra
aller, et au bas de la page une répartition qui les épuise toutes les neuf.

Restituées dans le vocabulaire d'aujourd'hui, les neuf points demandent :

\begin{enumerate}
\item ce qu'est un morphisme de \emph{structures asphériques}, et ce que
devient l'image d'une telle structure\footnote{Le mot que porte la page après
« morphisms of asph.\ structures ? » est illisible ; seul « images ? » se lit.
La restitution « image d'une structure » est celle de cette édition.} ;

\item ce qu'est un morphisme de \emph{contracteurs} ;

\item comment une structure s'induit sur une catégorie tranchée\footnote{La
page porte $M_{/a}$, d'une lecture douteuse : l'indice pourrait être une autre
lettre.} ;

\item pourquoi la théorie n'est pas auto-duale ;

\item à quelle condition un type de structure donne un modélisateur
canonique\footnote{L'adjectif qui qualifie « structure type » est abrégé en deux
lettres illisibles, et le mot « an » qui le précède est biffé. Rien sur la page
ne permet de le restituer.} ;

\item s'il existe des théorèmes d'existence de foncteurs test ;

\item quelles conditions de finitude on peut imposer dans un modélisateur
élémentaire\footnote{La phrase s'arrête sur l'adjectif : le substantif n'est pas
écrit. « modélisateur élémentaire » est la lecture de cette édition, appuyée sur
l'usage constant de l'expression ailleurs dans le même cercle de notes.} ;

\item si l'on sait calculer la cohomologie par un opérateur de bord ;

\item ce que donnent, comme modélisateurs de $(\mathrm{Cat})$, les variantes
finies et infinies des catégories test usuelles\footnote{Deux mots de ce point
sont illisibles sur la page, dont les deux substantifs que « infinite » et
« finite » qualifient. Ce qui est restitué ici est le squelette de la phrase, non
son contenu.}.
\end{enumerate}

La répartition du bas de page mérite d'être relevée, parce qu'elle partage
exactement les neuf points : 7 et 8 « may not be treated at all », 2, 3 et 9
iront dans la partie IV, 1, 4, 5 et 6 seront revus à Caen. C'est un plan de
travail, avec ses abandons.

\pagerange{2}{2}
\section*{Les catégories en jeu, et la question (page 2)}

Le deuxième feuillet est une liste de sept lignes et une question. Les six
premières lignes nomment les catégories test candidates, chacune avec sa
variante marquée d'un exposant $f$ :
\[
\Delta^{f}, \qquad \square,\ \square^{f}, \qquad \bigcirc,\ \bigcirc^{f},
\qquad \mathcal{B}_{f}(I).
\]
$\Delta$ est la catégorie des simplexes, $\square$ celle des cubes ;
l'exposant $f$ marque une variante restreinte, et la page ne dit pas
laquelle\footnote{La page ne donne aucune clef de ces signes. Que $\Delta$ soit
la catégorie des simplexes et $\square$ celle des cubes est sûr par l'usage ;
$\bigcirc$ et $\mathcal{B}_{f}(I)$ ne sont pas identifiés ici, et les identifier
serait deviner. De $\mathcal{B}_{f}(I)$ la page demande seulement si c'est une
catégorie test faible.}. En marge, une remarque sur $(\mathrm{Ord})$, la
catégorie des ensembles ordonnés : elle porte une structure asphérique, ou une
structure de contractilité.

Puis la question, seule sur sa ligne :
\[
\text{$0$-connexité totale} \ \Longrightarrow\ \text{asphéricité totale}\ ?
\]
La flèche n'a de pointe qu'à droite : c'est une implication, non une
équivalence. La page suivante répond, et répond mieux que la question ne
demandait — les deux conditions sont équivalentes.

\pagerange{3}{3}
\section*{L'asphéricité totale ne dépend pas du localisateur (page 3)}

\subsection*{Le vocabulaire}

Trois notions sont supposées connues sur cette page ; on les fixe ici.

Soit $A$ une petite catégorie et $\hat{A} = \mathrm{Fun}(A^{\circ},
\mathrm{Ens})$ la catégorie des préfaisceaux sur $A$.

Un \emph{localisateur fondamental} $W$ est une classe de foncteurs entre petites
catégories qu'on décide d'appeler équivalences : elle est faiblement saturée,
elle contient les foncteurs $A \to e$ pour toute $A$ ayant un objet final, et
elle vérifie une condition du type « théorème A » de Quillen. L'intérêt de
l'axiomatique est qu'il n'y en a pas qu'un : à chaque $W$ correspond une théorie
de l'homotopie, et une notion d'asphéricité.

Une petite catégorie $A$ est \emph{$W$-asphérique} si $A \to e$ appartient à
$W$ ; elle est \emph{totalement $W$-asphérique} si de plus le produit
$a \times b$ de deux représentables est $W$-asphérique dans $\hat{A}$, pour tous
objets $a, b$ de $A$.

Un objet $X$ de $\hat{A}$ est \emph{$0$-connexe} si $\pi_{0}(X)$ est réduit à un
point — autrement dit si la catégorie $A_{/X}$ est non vide et connexe. On note
$\hat{A}_{W_{0}}$ la sous-catégorie pleine des objets $0$-connexes.

\subsection*{La proposition}

\textbf{Proposition 4.} \emph{Soient $A$ une petite catégorie, $W$ un
localisateur fondamental, $W_{\infty}$ le localisateur fondamental le plus fin.
Les conditions suivantes sont équivalentes :}

\begin{enumerate}
\item[(i)] \emph{$A$ est totalement $W_{\infty}$-asphérique ;}
\item[(ii)] \emph{$A$ est totalement $W$-asphérique ;}
\item[(iii)] \emph{$A$ est non vide et, pour tous objets $a, b$ de $A$, le
produit $a \times b$ est $0$-connexe dans $\hat{A}$.}
\end{enumerate}

L'énoncé dit deux choses d'un coup. Que (i) et (ii) soient équivalentes, c'est
que l'asphéricité totale \emph{ne dépend pas} du localisateur choisi : la notion
qu'on croyait relative à $W$ est absolue. Et que (iii) leur soit équivalente,
c'est qu'elle ne dépend même pas de l'homotopie — la condition se lit sur les
seules composantes connexes.

\subsection*{La démonstration, et ce qu'elle suppose}

Le mécanisme tient en un encadrement. Tout localisateur fondamental $W$ vérifie
\[
W_{\infty} \ \subseteq\ W \ \subseteq\ W_{0},
\]
où $W_{0}$ est le localisateur fondamental le plus grossier. L'asphéricité
totale étant d'autant plus facile à obtenir que la classe des équivalences est
large, l'encadrement donne aussitôt
\[
\text{(i)} \Longrightarrow \text{(ii)} \Longrightarrow
\text{totalement } W_{0}\text{-asphérique},
\]
et il reste à voir que la dernière condition est exactement (iii), puis qu'elle
entraîne (i)\footnote{Sur la page les deux signes d'implication de cette ligne
manquent : la machine à écrire a ménagé les intervalles et rien n'y a été porté
à la main. Le sens est déterminé par la phrase qui suit — « it is immediate that
(iii) implies that $A$ is $0$-connected » — et par l'encadrement lui-même.}.

C'est immédiat, car $W_{0}$ est le localisateur pour lequel un foncteur est une
équivalence dès qu'il induit une bijection sur les composantes connexes :
$W_{0}$-asphérique et $0$-connexe sont le même mot\footnote{La page s'arrête sur
« namely » — elle allait nommer $W_{0}$, et la suite est à la page 351, qui n'est
pas dans la chemise. L'identification de $W_{0}$ à la classe des foncteurs
induisant une bijection sur $\pi_{0}$ est donc celle de cette édition, non celle
de la page ; c'est le localisateur fondamental maximal, et le seul candidat
compatible avec l'usage qui en est fait deux lignes plus haut.}.

Reste l'inclusion de gauche, et c'est là qu'il faut être précis sur ce que la
proposition suppose. Que les équivalences d'homotopie faibles usuelles forment
un localisateur fondamental est facile ; qu'elles forment le \emph{plus fin},
c'est-à-dire que $W_{\infty}$ soit contenu dans tout localisateur fondamental,
est une conjecture de Grothendieck. Elle a été démontrée par Denis-Charles
Cisinski une vingtaine d'années après ces feuillets\footnote{D.-C. Cisinski,
« Le localisateur fondamental minimal », \emph{Cahiers de topologie et géométrie
différentielle catégoriques} 45 (2004) ; voir aussi \emph{Les préfaisceaux comme
modèles des types d'homotopie}, Astérisque 308 (2006). La page de 1983 écrit
« $W_{\infty}$ the finest basic localizer (i.e.\ usual weak equivalence) » sans
plus de commentaire : le résultat y est en usage, non en démonstration. Rien
n'autorise à lire cette parenthèse comme autre chose que la conjecture prise
pour acquise dans un brouillon.}. La proposition 4, telle qu'elle est écrite,
en dépend.

\subsection*{Une remarque sur la pratique}

Le paragraphe qui précède la proposition n'est pas de la mathématique et c'est
peut-être ce que le feuillet a de plus rare. Grothendieck vient d'achever une
démonstration, s'interrompt, et écrit qu'après une minute de réflexion il se
trouve très bête d'avoir traîné pendant des mois la notion de catégorie
totalement $W$-asphérique sans avoir pris la peine de s'asseoir et de regarder
si elle dépendait vraiment de $W$ — et, du même coup, si la condition sur les
produits n'entraînait pas déjà le reste. « So let's clean it up in the long
last~! »

Ce qui est laissé sur la page, ce n'est pas le résultat mais le fait de ne pas
avoir posé plus tôt la question la plus simple. La proposition 4 tient en six
lignes ; le mois d'inattention qui la précède est écrit à côté d'elle.

\pagerange{4}{4}
\section*{La partie asphérique d'un modélisateur (page 4)}

Les deux derniers feuillets, en français, changent de sujet et de fortune.

\subsection*{La situation}

On se donne un modélisateur $M$, un foncteur $i : A \to M$ d'une petite
catégorie, et le foncteur qu'il induit
\[
i^{*} : M \longrightarrow \hat{A}, \qquad
i^{*}(x) = \mathrm{Hom}_{M}(i(-), x).
\]
On note $M_{as}$ la sous-catégorie pleine des objets asphériques de $M$. La
question de tout le feuillet est : \emph{quand $i^{*}$ est-il conservatif et
fidèle} — c'est-à-dire quand les objets de $A$ suffisent-ils à reconnaître les
objets de $M$ et à séparer les flèches ?

Le premier diagramme pose le cadre, avec un second modélisateur $M'$ au-dessus :

\begin{tikzcd}[column sep=large]
\hat{A} & M \arrow[l, "i^{*}"'] & M' \arrow[l, "f^{*}"']
\end{tikzcd}

et le second demande si l'on peut faire passer $i$ par la partie asphérique :

\begin{tikzcd}
A \arrow[r, "i"] \arrow[dr, dashed] & M\\
 & M_{as} \arrow[u]
\end{tikzcd}

\subsection*{Le « NB », et pourquoi il ne tient pas}

Grothendieck écrit alors, en gras, le cas modèle dont il espère tout tirer : si
$f : A \to B$ est un foncteur \emph{asphérique}, alors
$f^{*} : \hat{B} \to \hat{A}$ est conservatif et fidèle, « i.e.\ $f$ dominant »,
la dominance étant explicitée en marge :
\[
\forall\, b \in \mathrm{Ob}\,B,\ \exists\, a \in A \text{ et } f(a) \to b,
\qquad \text{c'est-à-dire } f_{/b} \neq \emptyset.
\]

Cet énoncé est faux, et il suffit de deux objets pour le voir\footnote{La lecture
retenue ici est celle des mots : « conservatif » et « fidèle » au sens ordinaire,
c'est-à-dire pour les isomorphismes et les égalités de flèches de préfaisceaux.
Si l'on entend « conservatif » au sens homotopique — $f^{*}$ reflète les
équivalences faibles — l'énoncé redevient vrai et c'est une propriété standard
des foncteurs asphériques. La page ne tranche pas ; mais sa dernière ligne, qui
demande que $M_{as}$ soit \emph{génératrice}, est une condition de séparation au
sens ordinaire, et c'est ce qui fait pencher cette édition. Il faut ajouter que
tout le travail qui entoure ce NB est raturé et repris deux fois : ce n'est pas
un énoncé arrêté qu'on corrige ici, c'est une note de travail dont on dit où elle
allait mal.}.

Soit $B$ la catégorie à deux objets $b_{0}, b_{1}$ et une seule flèche non
identique $\beta : b_{0} \to b_{1}$ ; soit $A = e$ la catégorie ponctuelle et
$f$ le foncteur d'image $b_{0}$.

\begin{tikzcd}[column sep=large]
e \arrow[r, "f"] & b_{0} \arrow[r, "\beta"] & b_{1}
\end{tikzcd}

Le foncteur $f$ est asphérique : $f_{/b_{0}}$ et $f_{/b_{1}}$ ont chacune un
seul objet et aucune flèche non identique, ce sont des catégories ponctuelles.
Il est donc en particulier dominant. Pourtant $f^{*} : \hat{B} \to \mathrm{Ens}$,
qui à un préfaisceau $X$ associe $X(b_{0})$, n'est ni conservatif ni fidèle :

\begin{itemize}
\item \emph{Non conservatif.} Soit $X$ le préfaisceau $X(b_{0}) = \{*\}$,
$X(b_{1}) = \emptyset$, et $Y$ le représentable $b_{1}$, de sorte que
$Y(b_{0}) = \{\beta\}$ et $Y(b_{1}) = \{\mathrm{id}\}$. Il existe un morphisme
$u : X \to Y$ — la condition de naturalité est vide en $b_{1}$ — et $f^{*}(u)$
est la bijection $\{*\} \to \{\beta\}$, alors que $u$ n'est pas un isomorphisme.

\item \emph{Non fidèle.} Soit $X$ le représentable $b_{1}$ et $Y$ le préfaisceau
$Y(b_{0}) = \{*\}$, $Y(b_{1}) = \{1, 2\}$. Les deux morphismes $X \to Y$ qui
envoient $\mathrm{id}_{b_{1}}$ sur $1$ et sur $2$ sont distincts, mais coïncident
en $b_{0}$ puisque $Y(b_{0})$ est ponctuel.
\end{itemize}

Rien de paradoxal : $B$ a un objet initial, donc $f$ est asphérique sans que
$\hat{B}$ ait la moindre raison de se laisser voir depuis un seul de ses objets.

\subsection*{Ce que la dominance caractérise exactement}

Deux lignes plus bas, la page affiche entre crochets un énoncé plus faible :
\[
\bigl[\ f^{*}(X) = \emptyset_{\hat{A}} \ \Longrightarrow\
X = \emptyset_{\hat{B}}\ \bigr].
\]
Celui-là est exact, et il est même équivalent à la dominance.

\emph{En effet.} Supposons $f_{/b} \neq \emptyset$ pour tout $b$, et soit $X$
avec $f^{*}(X) = \emptyset$, c'est-à-dire $X(f(a)) = \emptyset$ pour tout $a$.
Si $X(b)$ contenait un élément, la restriction le long d'une flèche
$f(a) \to b$ en produirait un dans $X(f(a))$ : contradiction. Réciproquement, si
un objet $b$ vérifie $f_{/b} = \emptyset$, le préfaisceau représenté par $b$ est
non vide et satisfait $f^{*}(b) = \mathrm{Hom}(f(a), b) = \emptyset$ pour tout
$a$. \qquad$\square$

La dominance ne dit donc rien de plus que ceci : $f^{*}$ reconnaît l'objet
initial. C'est très loin de la conservativité, et c'est exactement l'écart que
les ratures du feuillet sont en train de creuser.

\subsection*{Le bon critère}

Grothendieck y arrive, à la dernière ligne du feuillet, et cette fois c'est
juste : pour conclure que $i^{*}$ est fidèle, il faut que $M_{as}$ soit
\emph{génératrice}.

C'est en effet la caractérisation. Puisque $i^{*}(x)_{a} =
\mathrm{Hom}_{M}(i(a), x)$, deux flèches $u, v : x \to y$ de $M$ ont même image
par $i^{*}$ si et seulement si $u\varphi = v\varphi$ pour tout $a$ et tout
$\varphi : i(a) \to x$. Dire que $i^{*}$ est fidèle, c'est exactement dire que
la famille $\{i(a)\}_{a \in A}$ est \emph{séparatrice} dans $M$ ; et la
conservativité demande davantage encore\footnote{La conservativité de $i^{*}$
équivaut à ce que la famille soit \emph{fortement} génératrice, condition
strictement plus forte que la séparation, comme le montre déjà le
contre-exemple ci-dessus. La page ne distingue nulle part les deux, et ce qu'elle
appelle « génératrice » est ce qu'on appelle ici séparateur.}. C'est une
condition sur la \emph{taille} de la famille, non sur la nature de chaque
foncteur — ce que la dominance ne pouvait pas donner.

\pagerange{4}{5}
\section*{Deux contre-exemples (pages 4 et 5)}

Le critère une fois posé, la question devient : $M_{as}$ est-elle jamais
génératrice ? Grothendieck construit deux objets qui répondent non, et
franchement.

\subsection*{La partie asphérique peut être vide}

Prenons $A = \emptyset$. Alors $\hat{A}$ a exactement un objet — le préfaisceau
vide — et cet objet n'est pas $0$-connexe, puisque son $\pi_{0}$ est vide. Donc
\[
\hat{A} = \{e\}, \qquad \hat{A}_{as} = \emptyset.
\]
Rien à engendrer avec rien. « Contre-exemple ! », écrit la page.

\subsection*{La partie asphérique peut être ponctuelle sans rien engendrer}

Supposons maintenant $M_{as} = \{a\}$, réduite à un objet. Soit $A$ la
sous-catégorie pleine sur $a$ : c'est une catégorie à un seul objet, donc un
monoïde, celui de ses endomorphismes
\[
G = \underline{\mathrm{End}}(a),
\]
et $\hat{A}$ est la catégorie des $G$-ensembles à droite. Pour que les objets
asphériques de $\hat{A}$ soient les ensembles ponctuels, il faut que $G$ soit
trivial, et alors $\hat{A} \simeq (\mathrm{Ens})$ et $\hat{A}_{as}$ est la classe
des ensembles à un élément\footnote{La page affirme « donc nécessairement
$G = \{\mathrm{id}\}$ » sans un mot de justification, et la lecture des trois
premiers mots est douteuse. La raison est la proposition 4 elle-même, appliquée
à la catégorie à un objet : la condition (iii) exige que $a \times a$ soit
$0$-connexe, c'est-à-dire que l'action diagonale de $G$ sur $G \times G$ soit
transitive, ce qui pour un groupe force $|G| = 1$. Que les deux moitiés du
dossier se rejoignent ici n'est pas dit sur la page ; c'est la lecture de cette
édition.}. On obtient alors
\[
M_{as} = \{\, x \in \mathrm{Ob}\,M \ \mid \ \mathrm{Hom}(a, x)
\text{ est ponctuel} \,\}.
\]

Le contre-exemple tient en une ligne : soit $M$ une catégorie \emph{discrète} de
cardinal $\geq 2$. Aucun objet n'en reçoit de flèche d'un autre, la partie
asphérique ne peut donc rien séparer, et $i^{*}$ n'est pas fidèle. La condition
« $M_{as}$ génératrice » n'a rien d'automatique : elle est une hypothèse, et
c'est ce que ces deux exemples établissent.

\pagerange{5}{5}
\section*{La question laissée ouverte (page 5)}

Le feuillet finit sur une adjonction et une question.

Entre les deux modélisateurs, $f^{*} : M' \to M$ commute aux limites. Le
théorème du foncteur adjoint donne alors un adjoint à gauche
\[
f_{!} : M \longrightarrow M'
\]
dès que tout foncteur sur $M'$ commutant aux limites est représentable — ce qui
demande des hypothèses de solution-set, ou d'accessibilité, que la page se
contente de signaler comme « à vérifier »\footnote{« faire des hyp.\ d'\ldots\
accessibilité sur $f^{*}$ », dit la page, avec un mot illisible entre les deux.
Les noms modernes sont le théorème du foncteur adjoint de Freyd et les
catégories localement présentables de Gabriel--Ulmer, que Grothendieck pouvait
connaître ; la page ne les cite pas.}.

\begin{tikzcd}[column sep=large]
M \arrow[r, bend left=20, "f_{!}"] & M' \arrow[l, bend left=20, "f^{*}"]
\end{tikzcd}

Et la question, la dernière du dossier :
\begin{quote}
$f_{!}$ envoie-t-il $M_{as}$ dans $M'_{as}$ ?
\end{quote}

Autrement dit : l'adjoint à gauche respecte-t-il l'asphéricité ? Rien ne suit
sur le feuillet. C'est vers cette question que toute la page converge, et la
chemise s'arrête là\footnote{La notation du foncteur entre modélisateurs change
au fil des trois diagrammes du dossier — $f$ sur la page 4, $f^{*}$ sur la page
5 — sans qu'aucune ligne signale le changement. On a retenu partout $f^{*}$ pour
le foncteur $M' \to M$ et $f_{!}$ pour son adjoint à gauche, qui est l'état où
la page laisse les choses.}.

\end{document}
